% ============================================================
% ch04_implementation.tex
% WAL, Undo/Redo, Crash Simulation, Python Modules
% ============================================================

\chapter{Write-Ahead Logging Implementation}

\section{The WAL Protocol}

Every transaction in this system follows a strict 6-step sequence defined
in \texttt{transaction\_manager.py}. Steps must execute in this exact order:

\begin{table}[H]
\centering
\caption{WAL Transaction Steps}
\begin{tabularx}{\textwidth}{clX}
\toprule
\textbf{Step} & \textbf{Action} & \textbf{Details} \\
\midrule
1 & Read old\_balance & SELECT balance FROM Account WHERE account\_id = X \\
2 & \textbf{Write WAL log} & MongoDB insert (status=UNCOMMITTED) --- BEFORE any PG change \\
3 & Update PostgreSQL & UPDATE Account SET balance = new\_balance \\
4 & Insert Transaction\_Record & Records both old and new balance values \\
5 & conn.commit() & PostgreSQL transaction finalized \\
6 & mark\_committed() & MongoDB log updated to COMMITTED \\
\bottomrule
\end{tabularx}
\end{table}

\section{Deposit Implementation}

\begin{lstlisting}[style=pythonstyle, caption={Deposit method --- WAL sequence in Python}]
def deposit(self, account_number, amount, notes=""):
    conn   = get_pg_connection()
    cur    = conn.cursor()
    log_id = None
    try:
        # Step 1: Read old balance
        acc_id, old_bal = self._get_account(cur, account_number)
        new_bal = old_bal + amount

        # Step 2: Write WAL log BEFORE PostgreSQL update (WAL rule)
        log_id = self.logger.write_log(
            session_id  = self.session_id,
            account_id  = acc_id,
            txn_type    = "DEPOSIT",
            old_balance = old_bal,      # stored for UNDO
            new_balance = new_bal,      # stored for REDO
            amount      = amount,
        )

        # Step 3: Update PostgreSQL balance
        cur.execute("UPDATE Account SET balance = %s WHERE account_id = %s",
                    (new_bal, acc_id))

        # Step 4: Insert Transaction_Record
        self._record_transaction(cur, acc_id, "DEPOSIT",
                                 amount, old_bal, new_bal)

        # Step 5: COMMIT PostgreSQL
        conn.commit()

        # Step 6: Update WAL log to COMMITTED
        self.logger.mark_committed(log_id)

    except Exception as e:
        conn.rollback()                        # Reverse all PG changes
        if log_id:
            self.logger.mark_rolled_back(log_id, reason=str(e))
    finally:
        cur.close(); conn.close()
\end{lstlisting}

\section{Transfer --- Atomic Two-Leg Transaction}

A fund transfer involves two accounts and demonstrates ACID Atomicity:
both legs commit together or both are rolled back together.

\begin{lstlisting}[style=pythonstyle, caption={Transfer --- two WAL logs, one PG commit}]
# Step 2: Write BOTH WAL logs BEFORE any PostgreSQL update
log_out_id = self.logger.write_log(
    account_id=from_id, txn_type="TRANSFER_OUT",
    old_balance=from_bal, new_balance=from_new, amount=amount
)
log_in_id = self.logger.write_log(
    account_id=to_id, txn_type="TRANSFER_IN",
    old_balance=to_bal, new_balance=to_new, amount=amount
)

# Step 3: Apply both PostgreSQL updates
cur.execute("UPDATE Account SET balance=%s WHERE account_id=%s",
            (from_new, from_id))
cur.execute("UPDATE Account SET balance=%s WHERE account_id=%s",
            (to_new, to_id))

# Step 5: ONE commit covers both updates (Atomicity)
conn.commit()

# Step 6: Mark both WAL logs COMMITTED
self.logger.mark_committed(log_out_id)
self.logger.mark_committed(log_in_id)
\end{lstlisting}

% ────────────────────────────────────────────────────────────
\chapter{Crash Simulation}

\section{How a Crash is Simulated}

In \texttt{crash\_simulator.py}, a \texttt{RuntimeError} is raised between
Step 3 (PG update) and Step 5 (PG commit). This creates the same state
that a real OS crash would leave:

\begin{lstlisting}[style=pythonstyle, caption={crash\_simulator.py --- crash injection point}]
def simulate_crash(self, account_number, amount, txn_type="WITHDRAWAL"):
    conn   = get_pg_connection()
    cur    = conn.cursor()
    log_id = None
    try:
        acc_id, old_bal = _get_account(cur, account_number)
        new_bal = old_bal - amount

        # Step 2: WAL log written (UNCOMMITTED) --- this PERSISTS
        log_id = self.logger.write_log(
            session_id=self.session_id, account_id=acc_id,
            txn_type=txn_type, old_balance=old_bal,
            new_balance=new_bal, amount=amount,
            notes="CRASH_DEMO: will not commit"
        )

        # Step 3: PG balance changed
        cur.execute("UPDATE Account SET balance=%s WHERE account_id=%s",
                    (new_bal, acc_id))

        # *** CRASH HERE --- before conn.commit() ***
        raise RuntimeError("SIMULATED CRASH")

    except RuntimeError:
        conn.rollback()     # PG auto-rolls back (connection drop)
        # MongoDB log stays UNCOMMITTED <- recovery target
        return log_id
\end{lstlisting}

\section{State After Crash}

After \texttt{simulate\_crash()} returns:

\begin{table}[H]
\centering
\caption{Database State After Crash}
\begin{tabularx}{\textwidth}{lX}
\toprule
\textbf{Database} & \textbf{State} \\
\midrule
MongoDB wal\_logs & Log entry exists with \texttt{status = "UNCOMMITTED"} \\
PostgreSQL Account & Balance is restored to old\_balance (PG auto-rollback) \\
PostgreSQL Transaction\_Record & No record exists (INSERT was in the same transaction) \\
\bottomrule
\end{tabularx}
\end{table}

\begin{conceptbox}
\textbf{Important distinction for viva:}
In our simulation, PostgreSQL auto-rolls-back because we close the connection.
In a real crash, the OS terminates the process and PostgreSQL detects the
lost connection during its next startup. The MongoDB log remains UNCOMMITTED
in both cases --- that is the key state the Recovery Manager reads.
\end{conceptbox}

% ────────────────────────────────────────────────────────────
\chapter{Recovery Manager}

\section{Three-Phase Recovery}

The \texttt{RecoveryManager} implements a simplified ARIES-inspired process
in three sequential phases:

\begin{lstlisting}[style=pythonstyle, caption={recovery\_manager.py --- main entry point}]
def run_recovery(self):
    # Phase 1: Analysis
    uncommitted, committed = self._phase_analysis()

    # Phase 2: UNDO uncommitted transactions
    undo_count = self._phase_undo(uncommitted)

    # Phase 3: REDO committed transactions if PG state is wrong
    redo_count = self._phase_redo(committed)
\end{lstlisting}

\section{Phase 1: Analysis}

Reads all WAL logs from MongoDB and classifies them:

\begin{lstlisting}[style=pythonstyle, caption={Analysis phase}]
def _phase_analysis(self):
    uncommitted = self.logger.get_uncommitted()  # status="UNCOMMITTED"
    committed   = self.logger.get_committed()    # status="COMMITTED"
    return uncommitted, committed
\end{lstlisting}

\section{Phase 2: UNDO}

For each UNCOMMITTED log (newest first), restores \texttt{old\_balance}:

\begin{lstlisting}[style=pythonstyle, caption={UNDO phase --- restore old\_balance}]
def _phase_undo(self, uncommitted_logs):
    for log in uncommitted_logs:
        acc_id      = log["account_id"]
        old_balance = log["old_balance"]   # <- from WAL log
        log_id      = log["log_id"]

        # Read current PG balance
        cur.execute("SELECT balance FROM Account WHERE account_id=%s",
                    (acc_id,))
        pg_balance = float(cur.fetchone()[0])

        # Idempotency: skip if already restored
        if abs(pg_balance - old_balance) < 0.001:
            self.logger.mark_rolled_back(log_id, "Already restored")
            continue

        # Apply UNDO: restore old_balance in PostgreSQL
        cur.execute("UPDATE Account SET balance=%s WHERE account_id=%s",
                    (old_balance, acc_id))
        conn.commit()

        # Update WAL log -> ROLLED_BACK
        self.logger.mark_rolled_back(log_id, "UNDO by RecoveryManager")

        # Write recovery audit entry
        self.logger.write_recovery_log(
            action="UNDO", account_id=acc_id,
            before_bal=pg_balance, after_bal=old_balance,
            reason="CRASH_DETECTED"
        )
\end{lstlisting}

\section{Phase 3: REDO}

For each COMMITTED log, checks if PostgreSQL already reflects the committed
change. If not, re-applies \texttt{new\_balance}:

\begin{lstlisting}[style=pythonstyle, caption={REDO phase --- re-apply new\_balance if PG state is wrong}]
def _phase_redo(self, committed_logs):
    for log in committed_logs:
        acc_id      = log["account_id"]
        new_balance = log["new_balance"]   # <- from WAL log

        cur.execute("SELECT balance FROM Account WHERE account_id=%s",
                    (acc_id,))
        pg_balance = float(cur.fetchone()[0])

        # If PG already matches -> no REDO needed (idempotent)
        if abs(pg_balance - new_balance) < 0.001:
            continue

        # REDO: re-apply new_balance
        cur.execute("UPDATE Account SET balance=%s WHERE account_id=%s",
                    (new_balance, acc_id))
        conn.commit()

        self.logger.write_recovery_log(
            action="REDO", account_id=acc_id,
            before_bal=pg_balance, after_bal=new_balance,
            reason="COMMIT_LOSS"
        )
\end{lstlisting}

\section{Idempotency}

A key property of the recovery process is \textbf{idempotency} --- running
recovery multiple times produces the same result as running it once.

\begin{itemize}
  \item \textbf{UNDO idempotency}: If PG balance already equals
        \texttt{old\_balance}, UNDO is skipped.
  \item \textbf{REDO idempotency}: If PG balance already equals
        \texttt{new\_balance}, REDO is skipped.
\end{itemize}

This ensures the Recovery Manager is safe to run even if it is interrupted
partway through.

% ────────────────────────────────────────────────────────────
\chapter{Python Module Explanation}

\section{Module Overview}

\begin{table}[H]
\centering
\caption{Python Modules and ADBMS Concepts}
\begin{tabularx}{\textwidth}{llX}
\toprule
\textbf{Module} & \textbf{Concept} & \textbf{Why it exists} \\
\midrule
db\_config.py & Polyglot Persistence & Central DB connection management \\
mongo\_logger.py & WAL & All MongoDB log operations in one place \\
transaction\_manager.py & Transactions, ACID & Enforces WAL sequence for every transaction \\
crash\_simulator.py & Crash Recovery & Creates testable inconsistency scenarios \\
recovery\_manager.py & UNDO/REDO, ARIES & Detects and repairs post-crash inconsistency \\
main.py & --- & Interactive demo runner for viva/presentation \\
\bottomrule
\end{tabularx}
\end{table}

\section{db\_config.py}

\begin{lstlisting}[style=pythonstyle, caption={db\_config.py --- connection functions}]
def get_pg_connection():
    conn = psycopg2.connect(host="localhost", dbname="bank_recovery_db",
                            user="postgres", password="your_password")
    conn.autocommit = False   # Manual transaction control
    return conn

def get_mongo_db():
    client = MongoClient("mongodb://localhost:27017/")
    return client["bank_recovery_logs"]
\end{lstlisting}

\textbf{ADBMS Concept:} \texttt{autocommit=False} is the key setting.
It means every SQL statement is inside an implicit transaction until we
call \texttt{conn.commit()} or \texttt{conn.rollback()}.

\section{mongo\_logger.py}

The \texttt{write\_log()} method is the heart of the WAL system:

\begin{lstlisting}[style=pythonstyle, caption={MongoLogger.write\_log() --- WAL insert}]
def write_log(self, session_id, account_id, txn_type,
              old_balance, new_balance, amount, notes=""):
    log_id = self._next_log_id()
    self.col.insert_one({
        "log_id":      log_id,
        "session_id":  session_id,
        "account_id":  account_id,
        "txn_type":    txn_type,
        "old_balance": float(old_balance),   # UNDO uses this
        "new_balance": float(new_balance),   # REDO uses this
        "amount":      float(amount),
        "status":      "UNCOMMITTED",        # Always starts UNCOMMITTED
        "timestamp":   datetime.now(timezone.utc),
        "notes":       notes,
    })
    return log_id
\end{lstlisting}

\section{main.py --- Demo Menu}

\begin{lstlisting}[style=pythonstyle, caption={main.py menu structure}]
actions = {
    "1": menu_deposit,     # WAL + COMMIT demo
    "2": menu_withdraw,    # WAL + COMMIT or ROLLBACK
    "3": menu_transfer,    # Two-leg atomic transfer
    "4": menu_crash,       # Crash before COMMIT
    "5": menu_recovery,    # UNDO + REDO phases
    "6": menu_show_logs,   # View MongoDB WAL logs
    "7": menu_show_balances, # View PostgreSQL balances
}
\end{lstlisting}
